% 注：main.tex 中译注宏为 \znote / \znoteplain（任务书写作 \zhnote），此处按 main.tex 实际定义使用。
\chapter{Eichler--Shimura 关系}

% ===== PDF p.062 =====
\noindent{\small\emph{（本讲由 Ben Heuer 撰写）}}

\textbf{上次}：
\begin{itemize}
\item 有 $\GL_2(\Z)$-等变的双射
\[
\HH^\pm\to\Bigl\{\C\ \text{上的椭圆曲线}\ E\ \text{连同}\ \alpha\colon\Z^2\xrightarrow{\sim}H_1(E,\Z)\Bigr\}/\!\sim,
\]
\[
\tau\mapsto\bigl(E_\tau:=\C/(\Z+\Z\tau),\ \alpha\colon\Z^2\to\Z+\Z\tau,\ (1,0)\mapsto1,\ (0,1)\mapsto\tau\bigr).
\]
\item \textbf{推论}：
\[
\GL_2(\Z)\backslash\bigl(\GL_2(\widehat{\Z})\times\HH^\pm\bigr)\cong\Bigl\{\C\ \text{上的椭圆曲线}\ E\ \text{连同}\ \alpha\colon\widehat{\Z}^2\xrightarrow{\sim}T_E\Bigr\}/\!\sim,
\]
其中 $T_E:=\varprojlim_{N\in\N}E[N]$ 是 $E$ 的完整 adélic Tate 模。
\item 这可以定义在 $\Q$ 上！
\item 令 $\widehat{X}\to\Spec(\Q)$ 为表示函子
\[
\Q\text{-}\mathbf{Sch}\to\mathbf{Sets},\qquad S\mapsto\Bigl\{S\ \text{上的椭圆曲线}\ E\ \text{连同}\ \alpha\colon\widehat{\Z}^2\xrightarrow{\sim}T_E\Bigr\}/\!\sim
\]
的模概形。这里 $T_E=\varprojlim_N E[N]$ 是 $S$ 上有限平展群概形的逆向极限。
\item $-1\in\GL_2(\widehat{\Z})$ 在 $\widehat{X}$ 上作用平凡。
\item 对满足 $K\subseteq K_1(3)$ 的紧开子群 $K\subseteq\GL_2(\widehat{\Z})$（注意特别地 $-1\notin K$），有光滑拟射影曲线 $X_K=\widehat{X}/K$。
\item 例如：若 $K=K(N)$、$N\ge3$，则 $X_{K(N)}$ 参数化椭圆曲线 $E$ 连同平凡化 $(\Z/N)^2\xrightarrow{\sim}E[N]$。
\item 其他例子：$K=K_1(N)$ 给出 $\Z/N\hookrightarrow E[N]$，$K=K_0(N)$ 给出子群。
\end{itemize}

\textbf{下一个目标}：
\begin{itemize}
\item 用模曲线几何地重新解释 Hecke 算子。
\item 由模曲线模 $p$ 的几何，导出 Galois 作用与 Hecke 作用之间的基本关系：Eichler--Shimura 关系\idx{Eichler--Shimura 关系}。
\end{itemize}

\textbf{adélic 作用}：
\begin{itemize}
\item 显然，$\GL_2(\widehat{\Z})$ 通过 $\alpha\colon\widehat{\Z}^2\xrightarrow{\sim}T_E$ 上的预复合作用在 $\widehat{X}$ 上。
\item 但是：为了得到 Hecke 算子，我们需要 $\GL_2(\Af)$-作用 $\Rightarrow$ 要在准同源的意义下处理椭圆曲线。
\item 回忆：$\Af=\widehat{\Z}\otimes_\Z\Q$，且
\[
\GL_2(\Z)\backslash\bigl(\GL_2(\widehat{\Z})\times\HH^\pm\bigr)\cong\GL_2(\Q)\backslash\bigl(\GL_2(\Af)\times\HH^\pm\bigr).
\]
\item 在椭圆曲线上可以这样解释这一点。\znoteplain{为什么要「到准同源为止」：$\GL_2(\Af)$ 中 $p$-进方向的分量会改变 $T_E$ 这个格，因而在 $E$ 的同构意义下不能直接作用；但在准同源意义下，有理 Tate 模 $VE=T_E\otimes_\Z\Q$ 只依赖于 $E$ 的同源类，于是 $\GL_2(\Af)$ 得以作用。}
\item 回忆左边几何化为 $\widehat{X}$，它表示函子
\[
P\colon\Q\text{-}\mathbf{Sch}\to\mathbf{Sets},\qquad S\mapsto\{(E|S,\ \alpha\colon\widehat{\Z}^2\xrightarrow{\sim}T_E)\}/\text{同构}.
\]
\item 令 $VE=T_E\otimes_\Z\Q$ 为有理 adélic Tate 模，并定义
\[
P'\colon\Q\text{-}\mathbf{Sch}\to\mathbf{Sets},\qquad S\mapsto\{(E|S,\ \alpha\colon\Af^2\xrightarrow{\sim}VE)\}/\text{准同源}.
\]
\znoteplain{原文此处有脚注 14，但底本的提取文本中没有该脚注的正文，暂付阙如，待对照原书补入。}
\end{itemize}

% ===== PDF p.063 =====
\noindent\textbf{定义 8.1}\phantomsection\label{def:8.1}\quad 回忆：同源是指具有有限平坦核的满同态 $E\to E'$。\textbf{准同源}是指
\[
\{E\to E'\ \text{同源}\}\otimes_\Z\Q
\]
中的元素。\textbf{$p$-准同源}是指
\[
\{E\to E'\ \text{同源，其次数为}\ p\ \text{的幂}\}\otimes_\Z\Z\bigl[\tfrac1p\bigr]
\]
中的元素。
\begin{itemize}
\item 于是 $P\cong P'$。确实：
\begin{itemize}
\item 显然有自然变换 $P\to P'$。逆方向的略述如下：
\item 取 $(E,\alpha\colon\Af^2\xrightarrow{\sim}VE)\in P'(S)$。乘以同源 $N\colon E\to E$，直到 $\alpha^{-1}$ 限制到 $T_E$ 上给出 $T_E\to\widehat{\Z}^2$。把两边模 $M\in\N$ 约化。令 $D$ 为诱导映射
\[
E[M]\to(\Z/M)^2
\]
的核；当 $M\gg0$ 时它稳定下来。对偶同源 $E/D\to E$ 诱导出 $\alpha\colon\widehat{\Z}^2\xrightarrow{\sim}T(E/D)$。\znoteplain{这里「乘以同源 $N$」的作用是消去分母：$\alpha^{-1}(T_E)$ 是 $VE$ 中与 $\widehat{\Z}^2$ 可公度的格，乘 $N$ 倍之后可使它落进 $\widehat{\Z}^2$ 内（相差一个有限指数）；随后 mod $M$ 的下降再把这个有限指数商掉。}
\item 这在 $P(S)$ 中定义了唯一的代表元。于是 $P(S)\twoheadrightarrow P'(S)$。
\end{itemize}
\item 这表明自然的 $\GL_2(\widehat{\Z})$-作用延拓为 $\GL_2(\Af)$-作用。
\item 但这不固定到任何 $X_K$ 上的投影，因为它改变了 $E$。
\item \textbf{例}：$\begin{pmatrix}p&0\\0&1\end{pmatrix}$ 的作用是送 $E\mapsto E/D$，其中 $D\subseteq E[p]$ 是由 $\alpha(1,0)$ 在 $T_pE\twoheadrightarrow E[p]$ 下的像所生成的子群。
\item \textbf{事实}：这一切在紧化之后仍然成立。
\end{itemize}

\medskip
\noindent\textbf{Hecke 算子的几何解释}
\begin{itemize}
\item 令 $V=C^\infty(\GL_2(\Af)\times\HH^\pm)$ 为光滑函数
\[
f\colon\GL_2(\Af)\times\HH^\pm\to\C,\qquad (g,z)\mapsto f(g,z)
\]
构成的光滑 $\GL_2(\Af)$-表示，即：$f$ 被 $\GL_2(\Af)$ 的某个紧开子群固定，且对任意 $g\in\GL_2(\Af)$，$z\mapsto f(g,z)$ 是光滑的。
\item 固定一个素数 $p$。
\item 令 $\phi\in\cH(\GL_2(\Q_p))$ 为 $p$ 处 Hecke 代数中的任意元素，即具有紧支集的局部常值函数 $\phi\colon\GL_2(\Q_p)\to\C$。
\item 回忆对 $f\in V$ 有
\[
\phi*f(g,z)=\int_{\GL_2(\Q_p)}\phi(h)f(gh,z)\,dh,
\]
其中我们在 $\GL_2(\Q_p)$ 上取使 $K_p=\GL_2(\Z_p)$ 的体积为 $1$ 的 Haar 测度。
\item 上面的积分可以改写如下。
\item 考虑图表
\[
\begin{tikzcd}[column sep=3.2em, row sep=2.4em]
& \GL_2(\Af)\times\HH^{\pm}\times\GL_2(\Q_p)\arrow[dl,"q_1"']\arrow[dr,"q_2"] &\\
\GL_2(\Af)\times\HH^{\pm} & & \GL_2(\Af)\times\HH^{\pm}
\end{tikzcd}
\]
\end{itemize}

% ===== PDF p.064 =====
\[
q_1(g,z,h)=(g,z),\qquad q_2(g,z,h)=(gh,z).
\]
\begin{itemize}
\item 这样的图表称为\textbf{对应}（correspondence）。
\item 还有到第三个因子的投影：
\[
\pi\colon\GL_2(\Af)\times\HH^\pm\times\GL_2(\Q_p)\to\GL_2(\Q_p).
\]
\end{itemize}

\noindent\textbf{引理 8.2}\phantomsection\label{lem:8.2}\quad 对 $f\colon\GL_2(\Af)\times\HH^\pm\to\C$（$f\in V$），有
\[
\phi*f=q_{1,!}\bigl(q_2^*(f)\cdot\pi^*(\phi)\bigr),
\]
其中 $(-)^*$ 表示函数的拉回，$(-)_!$ 表示沿纤维积分（由于 $\phi$ 有紧支集，这是有意义的）。

\noindent\textbf{证明}．$q_1$ 在 $(g,z)\in\GL_2(\Af)\times\HH^\pm$ 上的纤维是
\[
q_1^{-1}(g,z)=\{(g,z,h)\mid h\in\GL_2(\Q_p)\}\cong\GL_2(\Q_p),
\]
而 $q_2^*(f)\cdot\pi^*(\phi)$ 是函数 $(g,z,h)\mapsto f(gh,z)\phi(h)$。施加 $q_{1,!}$，即沿上述纤维积分，得到
\[
\int_{\GL_2(\Q_p)}\phi(h)f(gh,z)\,dh=\phi*f(g,z).
\]
注意这是良定义的，因为 $\phi$ 有紧支集。\hfill$\square$

\medskip
\noindent\textbf{过渡到有限级}
\begin{itemize}
\item 设 $N$ 与 $p$ 互素，并令 $K:=K_1(N)\subseteq\GL_2(\Af)$。则 $K_p=\GL_2(\Z_p)$。
\item 记 $K^p$ 为 $K$ 的 $p$-外部分，即 $K=K^pK_p$ 且 $K^p\cap K_p=1$。
\item 假设 $f$ 被 $K$ 固定，即 $f$ 由某个函数 $\tilde f\colon\GL_2(\Af)/K\times\HH^\pm\to\C$ 拉回得到。
\item 假设 $\phi\in\cH(\GL_2(\Q_p),K_p)$，即 $\phi$ 是 $K_p$-双不变的。
\item 于是 $\phi$ 是某个具有有限支集的函数 $\tilde\phi\colon K_p\backslash\GL_2(\Q_p)/K_p\to\C$ 的拉回。
\item 此时 $\phi*f$ 仍是 $K$-不变的，即它是某个函数 $\widetilde{\phi*f}\colon\GL_2(\Af)/K\times\HH^\pm\to\C$ 的拉回。
\item 于是在保持记号不变的情况下，可以把上面的图表改为
\[
\begin{tikzcd}[column sep=2.2em, row sep=2.8em]
& \bigl(\GL_2(\Af)/K^p\times\HH^{\pm}\bigr)\times_{K_p}\GL_2(\Q_p)/K_p\arrow[dl,"q_1"']\arrow[dr,"q_2"] &\\
\GL_2(\Af)/K\times\HH^{\pm} & & \GL_2(\Af)/K\times\HH^{\pm}
\end{tikzcd}
\]
这里 $-\times_{K_p}-$ 是第 3 节的收缩积（即对作用 $k\cdot(g,z,h):=(gk^{-1},z,kh)$ 取商），且
\[
q_2([g,z,h])=[gh,z],\qquad q_1([g,z,h])=[g,z],
\]
其中方括号表示等价类。
\end{itemize}

% ===== PDF p.065 =====
\begin{itemize}
\item 我们再次有投影
\[
\pi\colon\bigl(\GL_2(\Af)/K^p\times\HH^\pm\bigr)\times_{K_p}\GL_2(\Q_p)/K_p\longrightarrow K_p\backslash\GL_2(\Q_p)/K_p,\qquad [g,z,h]\mapsto[h].
\]
\item 接下来，在左边对 $\GL_2(\Q)$ 取商：回忆这给出复流形
\[
X_{\Gamma_1(N)}(\C)=\GL_2(\Q)\backslash\bigl(\GL_2(\Af)/K\times\HH^\pm\bigr).
\]
\item 然后需要把函数 $f$ 换成该流形上 $\omega_k$（$k\in\Z$）的截面。
\item 更精确地说，我们得到
\begin{equation}\tag{1}\label{eq:8.1}
\begin{tikzcd}[column sep=3.2em, row sep=2.8em]
& \begin{gathered}\GL_2(\Q)\backslash\bigl(\GL_2(\Af)/K^p\times\HH^{\pm}\bigr)\\\times_{K_p}\GL_2(\Q_p)/K_p\end{gathered}\arrow[dl,"q_1"']\arrow[dr,"q_2"] &\\
X_{\Gamma_1(N)}(\C) & & X_{\Gamma_1(N)}(\C)
\end{tikzcd}
\end{equation}
\item 并且对 $f\in H^0(X_{\Gamma_1(N)}(\C),\omega_k)$，即级为 $K_1(N)$ 的弱模形式，有
\[
\phi*f=q_{1,!}\bigl(q_2^*(f)\cdot\pi^*(\phi)\bigr).
\]
\item 为使之精确，要用到：有典范同构
\[
q_2^*(\omega_k)\cong q_1^*(\omega_k),
\]
因为 $\omega_k$ 是由 $\HH^\pm$ 上一个 $\GL_2(\Q)$-等变线丛的拉回来定义的。
\item 于是可以把 $q_2^*(f)\cdot\pi^*(\phi)\in q_2^*(\omega_k)$ 看作 $q_1^*(\omega_k)$ 的截面。
\item 于是可以沿纤维求和，得到截面
\[
q_{1,!}\bigl(q_2^*(f)\cdot\pi^*(\phi)\bigr)\in\Gamma\bigl(X_{\Gamma_1(N)}(\C),\omega_k\bigr).
\]
\item 这是模形式上 Hecke 作用的一个更几何的解释。
\end{itemize}

\medskip
\noindent\textbf{用模来解释}
\begin{itemize}
\item 在椭圆曲线上，式 (1) 对应一个 \textbf{Hecke 对应}：
\[
\begin{tikzcd}[column sep=4em, row sep=2.4em]
& \{(\iota\colon E_1\dashrightarrow E_2,\ \alpha)\}\arrow[dl,"q_1"']\arrow[dr,"q_2"] &\\
\{(E_1,\alpha)\} & & \{(E_2,\alpha')\}
\end{tikzcd}
\]
其中
\begin{itemize}
\item $(E_1,\alpha)$ 是 $\C$ 上的椭圆曲线带 $\Gamma_1(N)$-级结构 $\alpha\colon\Z/N\Z\hookrightarrow E_1[N]$；
\item 等价地，这是带一个恰为 $N$ 阶的点 $Q:=\alpha(1)\in E_1[N](\C)$ 的数据；
\item $(\iota\colon E_1\dashrightarrow E_2,\alpha)$ 是 $\C$ 上椭圆曲线 $E_1,E_2$ 连同 $E_1$ 上的 $\Gamma_1(N)$-级结构 $\alpha$ 以及一个 $p$-准同源 $\iota$ 的数据；
\item $q_1\colon(\iota\colon E_1\dashrightarrow E_2,\alpha)\mapsto(E_1,\alpha)$；
\item $q_2\colon(\iota\colon E_1\dashrightarrow E_2,\alpha)\mapsto(E_2,\alpha')$，其中 $\alpha'$ 是复合
\[
\Z/N\Z\xrightarrow{\ \alpha\ } E_1[N]\xrightarrow{\ \iota\ } E_2[N]
\]
（用到：由于 $(p,N)=1$，$\iota$ 在 $N$-挠部分上是同构）。等价地，它把 $Q$ 送到 $E_2[N]$ 中的像 $Q'$。\znoteplain{对应 $q_1,q_2$ 的用法：$q_2^*(f)\pi^*(\phi)$ 是中间空间上「带权」的函数，沿 $q_1$ 的纤维积分（纤维有限时就是求和）把它压回 $X_{\Gamma_1(N)}$；这正是 Hecke 算子 $\sum_{D}[E/D,Q+D]$ 的几何写法。}
\end{itemize}
\end{itemize}

% ===== PDF p.066 =====
\begin{itemize}
\item 映射 $\pi$ 把 $(\iota\colon E_1\dashrightarrow E_2,\alpha)$ 送到同构
\[
T_p(\iota)\colon T_pE_1\otimes_{\Z_p}\Q_p\xrightarrow{\sim}T_pE_2\otimes_{\Z_p}\Q_p
\]
的类。更具体地，任取 $T_pE_1$ 与 $T_pE_2$ 的基，则 $T_p(\iota)$ 由矩阵 $A\in\GL_2(\Q_p)$ 表示。$A$ 在 $K_p\backslash\GL_2(\Q_p)/K_p$ 中的类与基的选取无关。
\item \textbf{问题}：$p$-准同源构成的函子不可表（只能由归纳概形表示）。
\item 然而：我们主要关心由下列双陪集的特征函数给出的 $\phi$：
\[
K_p\begin{pmatrix}p&0\\0&p\end{pmatrix}K_p\ \longleftrightarrow\ \tilde S_p,\qquad
K_p\begin{pmatrix}p&0\\0&1\end{pmatrix}K_p\ \longleftrightarrow\ \tilde T_p .
\]
\item 这时定义积分中的项 $\phi(h)$ 的支集就在这个双陪集上。$\rightsquigarrow$ 于是 Hecke 对应的支集落在可由概形表示的子空间上！
\item 对 $\tilde S_p$，考虑形如
\[
([p]\colon E\to E,\ \alpha)
\]
的组 $(\iota\colon E_1\dashrightarrow E_2,\alpha)$，即 $E_1=E_2$ 且 $\iota\colon E_1\to E_2$ 是乘以 $[p]$。
\item 对 $\tilde T_p$，则需要考虑
\[
(\iota\colon E_1\to E_2,\ \alpha),
\]
其中 $\iota\colon E_1\to E_2$ 是次数为 $p$ 的同源。
\item 等价地，这是 $E_1$ 的一个秩为 $p$ 的循环子群 $\ker\iota$ 的数据。
\item \textbf{上次}：这些现在都可由概形表示，甚至可以在 $\Spec(\Q)$ 或适当的 $\Spec(\Z[1/N])$ 上表示！
\end{itemize}

\medskip
\noindent\textbf{模曲线上的 Hecke 作用——代数几何的观点}
\begin{itemize}
\item 现在我们过渡到有限级 $K=K_1(N)$，其中 $p\nmid N$。
\item 相应的模曲线 $X_{\Gamma_1(N)}$（在 $\Z[\frac1N]$ 上）表示 $\Z[\frac1N]$-概形 $S$ 上的对 $(E,Q)$：$S$ 上的椭圆曲线 $E$ 连同点 $Q\in E[N](S)$。
\item 我们从 $\langle p\rangle:=\tilde S_p$ 开始。由上述，只需把 $Q$ 乘以 $p$。
\item 为此，利用 $(\Z/N\Z)^\times$ 在 $X_{\Gamma_1(N)}$ 上的作用。
\item 由上述，$\tilde S_p$ 的对应现在成为
\[
\begin{tikzcd}[column sep=4em, row sep=2.4em]
& X_{\Gamma_1(N)}\arrow[dl,"\mathrm{id}"']\arrow[dr,"p\,\cdot"] &\\
X_{\Gamma_1(N)} & & X_{\Gamma_1(N)}
\end{tikzcd}
\]
\item \textbf{事实}：对任意 $a\in(\Z/N\Z)^\times$，有典范同构 $a^*\omega=\omega$。
\item 由于右边的映射是同构，沿纤维的积分是平凡的。
\item 一切都可以延拓到紧化上。
\item 综合上述，我们证明：
\end{itemize}

\noindent\textbf{命题 8.3}\phantomsection\label{prop:8.3}\quad 考虑由复合
\[
M_k(\Gamma_1(N),\Z[\tfrac1N])=H^0(X^*_{\Gamma_1(N)},\omega^{\otimes k})\xrightarrow{\ p^*\ }H^0(X^*_{\Gamma_1(N)},p^*\omega^{\otimes k})=M_k(\Gamma_1(N),\Z[\tfrac1N])
\]
定义的算子。则它基变换到 $\C$ 之后就是 Hecke 算子 $\tilde S_p$。

% ===== PDF p.067 =====
\begin{itemize}
\item 对 $\tilde T_p$，代替同源 $(\iota\colon E\to E',\alpha)$，我们参数化三元组 $(E,D,Q)$，其中 $D=\ker\iota$ 是秩为 $p$ 的有限循环子群。得到模函子的映射
\[
\pi_1\colon(E,D,Q)\mapsto(E,Q),\qquad \pi_2\colon(E,D,Q)\mapsto(E/D,Q'),
\]
\[
Q':=Q+D=\text{$Q$ 在}\ E[N]\to(E/D)[N]\ \text{下的像}.
\]
\item 它们诱导出 Hecke 对应
\[
\begin{tikzcd}[column sep=4em, row sep=2.4em]
& X_{\Gamma_0(p)\cap\Gamma_1(N)}\arrow[dl,"\pi_1"']\arrow[dr,"\pi_2"] &\\
X_{\Gamma_1(N)} & & X_{\Gamma_1(N)}
\end{tikzcd}
\]
\item \textbf{事实}：$\pi_1$ 与 $\pi_2$ 延拓到紧化 $X^*_{\Gamma_0(p)\cap\Gamma_1(N)}\to X^*_{\Gamma_1(N)}$。
\item \textbf{事实}：$\pi_1$ 与 $\pi_2$ 是次数为 $p+1$ 的有限平坦态射，在紧化上亦然。
\item \textbf{事实}：有典范同构 $\pi_1^*\omega=\pi_2^*\omega$（在 $\Spec(\Q)$ 上）。
\item 由于 $\pi_1$ 的纤维有限，Hecke 算子定义中出现的积分变成离散和，可以解释为迹映射
\[
\Tr_\pi\colon\pi_{1,!}\pi_1^*\omega\to\omega .
\]
\item 综合上述，我们证明：
\end{itemize}

\noindent\textbf{命题 8.4}\phantomsection\label{prop:8.4}\quad 把算子
\[
H^0(X^*_{\Gamma_1(N)},\omega^{\otimes k})\xrightarrow{\ \pi_2^*\ }H^0(X^*_{\Gamma_1(N)\cap\Gamma_0(p)},\omega^{\otimes k})\xrightarrow{\ \Tr_{\pi_1}\ }H^0(X^*_{\Gamma_1(N)},\omega^{\otimes k})
\]
基变换到 $\C$ 之后就是 Hecke 算子 $\tilde T_p$。

\begin{itemize}
\item \textbf{注}：这一切已经定义在 $\Q$ 上（甚至 $\Z[\frac1N]$ 上）！于是有
\end{itemize}

\noindent\textbf{推论 8.5}\phantomsection\label{cor:8.5}\quad $\tilde S_p,\tilde T_p$ 在 $M_k(\Gamma_1(N),\C)$ 上的特征值是代数的。

\begin{itemize}
\item 为给出 $\tilde T_p$ 更显式的描述，用除子重新解释。
\item 令 $\Div(X_{\Gamma_1(N)})$ 为 $X$ 上的除子集合（$\approx$ 点的形式和）。
\item 于是可以把 $\langle p\rangle:=\tilde S_p$ 与 $\tilde T_p$ 重新解释为 $\Div(X_{\Gamma_1(N)})$ 上的算子：
\[
\tilde S_p\colon\Div(X_{\Gamma_1(N)})\to\Div(X_{\Gamma_1(N)}),\qquad [E,Q]\mapsto[E,p\cdot Q],
\]
\[
\tilde T_p\colon\Div(X_{\Gamma_1(N)})\to\Div(X_{\Gamma_1(N)}),\qquad [E,Q]\mapsto\sum_{D\subseteq E[p]}[E/D,Q+D],
\]
其中 $E$ 是椭圆曲线，$Q\in E[N]$ 是 $N$ 阶的点。
\end{itemize}

\medskip
\noindent\textbf{特征 $p$ 中的模曲线}
\begin{itemize}
\item 设 $E$ 是特征 $p$ 的概形 $S$ 上的椭圆曲线。
\item 则当 $p\nmid N$ 时，$E[N]$ 是平展的。
\item 态射 $[p]\colon E\to E$ 分解为
\[
E\xrightarrow{\ F\ }E^{(p)}\xrightarrow{\ V\ }E,
\]
即 Frobenius 同源与 Verschiebung 同源。这里 $E^{(p)}=E\times_{S,F}S$。
\item 于是总有 $\ker F\subseteq E[p]$。它是连通的，即 $\ker F(S)=0$。
\end{itemize}

\noindent\textbf{定义 8.6}\phantomsection\label{def:8.6}\quad 在 $S=\F_p$ 上，有两种情形：

% ===== PDF p.068 =====
\begin{enumerate}
\item[(1)] $E[p](\Fq)\cong\Z/p\Z$。此时 $E$ 称为\textbf{普通的}（ordinary），且 $E[p]=\mu_p\times\Z/p\Z$。
\item[(2)] $E[p](\Fq)=0$。此时 $E$ 称为\textbf{超奇异的}（supersingular），且 $E[p]$ 连通。
\end{enumerate}
\znote{原文此定义（1）(2) 两条作 $E(\F_p)=\Z/p\Z$、$E(\F_p)=0$，按 $E[p](\Fq)$ 译：普通时 $E[p]\cong\mu_p\times\Z/p\Z$，超奇异时 $E[p]$ 连通。[来源：Goren, \emph{Hilbert Modular Forms}, Example 3.10，\url{https://www.wstein.org/people/goren/hilbert_modular_forms_book.pdf}（2026-11-26 访问）。]}

\begin{itemize}
\item 「超奇异」说的并不是光滑性。它只意味着「非常特殊」：
\item \textbf{事实}：$\Fq$ 上超奇异椭圆曲线只有有限多个同构类。其余都是普通的。
\item 回忆：$X_{\Gamma_1(N)}$ 定义在 $\Z[\frac1N]$ 上。
\item 因此可以约化到 $\F_p$，得到模曲线 $X_{\Gamma_1(N),\F_p}\to\Spec(\F_p)$。
\item 类似地，$X_{\Gamma_1(N)\cap\Gamma_0(p)}$ 表示 $(E,D,Q)$，其中 $D\subseteq E[p]$ 秩为 $p$。
\item Deligne--Rapoport：纤维 $X_{\Gamma_1(N)\cap\Gamma_0(p),\F_p}\to\Spec(\F_p)$ 具有如下形状
\[
\begin{tikzcd}[column sep=4.6em, row sep=2.6em]
& X_{\Gamma_1(N)\cap\Gamma_0(p),\F_p}\arrow[dl,"\mathrm{id}"']\arrow[dr,"F"] &\\
X_{\Gamma_1(N),\F_p}\ (D=\mu_p) & & X_{\Gamma_1(N),\F_p}\ (D=\Z/p\Z)
\end{tikzcd}
\]
\item 即两个 $X_{\Gamma_1(N),\F_p}$ 的拷贝，分别对应 $D=\mu_p$ 与 $D=\Z/p\Z$ 作为 $E[p]$ 的子群，它们在超奇异点处横截相交（图中记为黑点）。
\item 投影在一个拷贝上是恒等（$\deg=1$），在另一个上是 Frobenius（$\deg=p$）。
\end{itemize}

\medskip
\noindent\textbf{Eichler--Shimura 关系}

\noindent [Diamond--Shurman]，[Conrad: Appendix to Serre's...]
\begin{itemize}
\item Eichler--Shimura 关系把 $\tilde T_p$ 模 $p$ 的约化用 Frobenius 表达出来：
\item 基变换到代数闭域 $\Fq$。
\item 考虑态射 $\Div^0(X^*_{\Gamma_1(N),\Fq})\to\Pic^0(X^*_{\Gamma_1(N),\Fq})$。
\item \textbf{事实}：$\Pic^0(X^*_{\Gamma_1(N),\Fq})$ 是 $\Fq$ 上某个群概形的点。特别地，乘以 $p$ 通过 Frobenius $F$ 分解：存在 $V=F^\vee$ 使得 $p=V\circ F=F\circ V$。
\end{itemize}

\noindent\textbf{定理}（Eichler--Shimura 关系）\quad 在 $\Pic^0(X^*_{\Gamma_1(N),\Fq})$ 中，有
\[
\tilde T_p=F+\langle p\rangle V .
\]
\begin{itemize}
\item \textbf{注}：对 $X^*_{\Gamma_0(N),\Fq}$ 也有一个版本，众所周知它写作
\[
\tilde T_p=F+V .
\]
把上式乘以 $F$ 并用 $FF^\vee=p$，得到第二个重要的变体
\[
F^2-\tilde T_pF+\langle p\rangle p=0 .
\]
\znoteplain{这一变体的来历：由 $\tilde T_p=F+\langle p\rangle V$ 两边右乘 $F$，得 $\tilde T_pF=F^2+\langle p\rangle VF$；再用 $VF=p$ 即得上式。第 9 讲构造 Galois 表示时，特征多项式 $X^2-a_pX+\chi(p)p$ 正来源于此。}
\end{itemize}

\noindent\textbf{证明}（略述）．设 $E$ 是 $\breve{\Z}_p=W(\Fq)$ 上的普通椭圆曲线。令 $C\subseteq E[p]$ 为「典范子群」：它由
\[
E[p](\breve{\Z}_p)\to E[p](\Fq)
\]
的核生成。

% ===== PDF p.069 =====
\begin{itemize}
\item \textbf{关键事实}：对秩为 $p$ 的循环子群概形 $D\subseteq E[p]$，
\begin{enumerate}
\item[(1)] 同源 $E\to E/D$ 约化为 $F\colon E\to E^{(p)}$ $\iff$ $D=C$；
\item[(2)] 同源 $E\to E/D$ 约化为 $V\colon E\to E^{(p-1)}$ $\iff$ $D\ne C$。
\end{enumerate}
\item 这里 $E^{(p-1)}:=E\times_{\F_p,F^{-1}}\F_p$，即沿 Frobenius 的逆作基变换。
\item $F$ 把 $Q\in E[N]$ 送到基变换 $Q^{(p)}\in E^{(p)}[N]=E[N]^{(p)}$。
\item 由于 $V\circ F=p$ 作用在 $E[N]$ 上，这蕴含：$V$ 把 $Q$ 送到 $pQ^{(p-1)}\in E^{(p-1)}[N]$。
\item 回忆：$E[p]$ 在 $W(\Fq)$ 上有 $p+1$ 个不同的子群 $D$。
\item 因此 $\tilde T_p([E,Q])=\sum_{D\subseteq E[p]}[E/D,Q+D]$（在 $\Div(X_{\Gamma_1(N)})$ 中）约化为
\[
[E^{(p)},Q^{(p)}]+p[E^{(p-1)},pQ^{(p-1)}]
\]
（在 $\Div(X_{\Gamma_1(N),\Fq})$ 中）。
\item 第一个求和项是 $F[E,Q]$。
\item 按定义，$\langle p\rangle[E^{(p-1)},Q^{(p-1)}]=[E^{(p-1)},pQ^{(p-1)}]$。
\item 现在过渡到 $\Pic^0(X^*_{\Gamma_1(N),\Fq})$。这里 $p=VF$。于是有
\[
p[E^{(p-1)},Q^{(p-1)}]=pF^{-1}[E,Q]=V[E,Q].
\]
（严格地说，需要取 $[E,Q]-[E',Q']$ 才能落在 $\Div^0$ 中。）\hfill$\square$
\end{itemize}
